\label{sec:results}

\subsection{Experimental Setup}

We evaluate M.1 economic character weights against the geographic baseline
on three states.
LODES WAC 2020 data was downloaded and aggregated to census tracts using
\texttt{bisect fetch --type lodes --year 2020}.
All experiments use:
\begin{itemize}
  \item Algorithm: \texttt{standard-bisect} structure
  \item Seed: 42 (single seed)
  \item Blend factor: $\alpha = 0.5$ (\texttt{--econ-alpha 0.5})
  \item Year: 2020 decennial census population
\end{itemize}

\subsection{Proportionality Gap Results}

Table~\ref{tab:main-results} reports the proportionality gap (D-seat share
minus D-vote share) under geographic and economic character weights.

\begin{table}[ht]
\centering
\begin{threeparttable}
\caption{Proportionality gap: geographic vs.\ economic character weights.
  More negative = more Republican-favoring.
  $\Delta$ Gap is economic-character minus geographic (positive = improved).}
\label{tab:main-results}
\begin{tabular}{lcrrrr}
  \toprule
  State & $k$ & D vote (\%) & Geographic & Economic-char & $\Delta$ Gap \\
  \midrule
  North Carolina & 14 & 49.3 & $-$20.7\,pp & $-$6.4\,pp & $+$14.3\,pp \\
  Wisconsin\tnote{$\dagger$}      &  8 & 50.3 & $-$12.8\,pp & $-$12.8\,pp & 0.0\,pp \\
  Texas\tnote{$\dagger$}          & 38 & 47.2 & $-$31.4\,pp & $-$34.0\,pp & $-$2.6\,pp \\
  \bottomrule
\end{tabular}
\begin{tablenotes}
  \small
  \item[$\dagger$] Single-run result (seed=42). Multi-seed variance not yet measured.
\end{tablenotes}
\end{threeparttable}
\end{table}

\subsection{NC-14: The Research Triangle and Charlotte Corridor}

North Carolina provides the clearest demonstration that economic character
can serve as a meaningful community signal.
The 14.3\,pp gap improvement reflects two concentrated economic geography
features that act as natural cut seams under M.1 weights:

\paragraph{Research Triangle Park.}
The RTP complex---spanning tracts in Durham, Orange, and Wake counties---is one
of the largest planned research and development parks in the world.
Its census tracts have employment intensity (JPR) values well above 0.50,
reflecting tens of thousands of pharmaceutical, biotechnology, and information
technology jobs concentrated in facilities that employ far more workers than
live on-site.
Under geographic weights, these high-JPR tracts are bundled with their
residential neighbors: the geographic boundary length between RTP tracts
and adjacent Cary or Morrisville residential tracts is long (they share
extensive perimeter), so METIS treats them as expensive to cut.
Under M.1 economic character weights, the JPR contrast ($\mathrm{JPR}_{\mathrm{RTP}} \approx 0.7$
vs.\ $\mathrm{JPR}_{\mathrm{suburb}} \approx 0.03$) produces very low cosine
similarity between these adjacent tracts (approximately 0.23), making the cut cheaper.
The resulting partition separates RTP from its residential neighbors more cleanly,
producing districts that either contain the employment cluster or the residential
neighborhoods surrounding it---not both.

\paragraph{Charlotte commercial corridor.}
The I-85 and South Boulevard commercial corridors in Mecklenburg County
have high commercial intensity ($\mathrm{CI} > 0.4$) from their concentration
of retail, financial services, and professional services employers.
Under M.1 weights, the commercial-to-residential boundary becomes the preferred
cut seam, separating the commercial corridor from the residential districts
surrounding it.

The net effect of these two separations is 2 additional Democratic districts
(6 vs.\ 4), as employment centers that were previously bundled within
Republican-leaning residential districts become either their own districts
or are assigned to adjacent urban districts.

\subsection{WI-8: No Effect}

Wisconsin's null result ($\Delta = 0$\,pp) is consistent with the M-track
prediction for states with diffuse economic geography.
Wisconsin has several medium-sized cities (Milwaukee, Madison, Green Bay,
Appleton, Racine) but lacks the sharp employment concentration contrasts
of NC.
Its commercial and light industrial zones are spread across many small parcels
interspersed with residential development, rather than concentrated in large
distinct employment clusters.
The cosine similarities between adjacent WI tracts are more uniformly
distributed---there are fewer pairs with very low similarity that would
create strong cut seams.
The M.1 modifier applies, but it modulates all edges roughly proportionally,
producing no material change in the partition.

This null result confirms that M.1 is working correctly: a communities-of-interest
signal should have no effect when there are no concentrated communities to preserve.

\subsection{TX-38: Slight Worsening}

Texas shows a 2.6\,pp worsening (from $-$31.4\,pp to $-$34.0\,pp), a
directional result consistent with the economic-partisan misalignment hypothesis.
Houston's petrochemical refining and logistics zones and Dallas-Fort Worth's
warehousing and distribution corridors have high industrial fraction values.
These industrial zones are located in suburban and exurban Republican-leaning
areas.
Under M.1 weights, the industrial-to-residential boundaries in these areas
become preferred cut seams, which separates Republican-leaning industrial
suburbs from adjacent areas in a way that produces marginally fewer Democratic
districts.

The effect is modest ($-$2.6\,pp) and within the expected single-seed variance
range.
It should be confirmed with multi-seed runs before treating it as a reliable
finding.

\subsection{Within-District Economic Variance: Phase 2}

The M-track primary metric $\overline{\sigma}_{\mathrm{JPR}}$ (mean within-district
standard deviation of jobs per resident, equation~\ref{eq:within-var}) was not
computed in the current pipeline run.
Computing this metric requires per-district tract-level output with JPR values,
which requires a post-processing analysis step that is designated Phase~2.

The qualitative prediction is clear from the NC mechanism:
when employment centers like RTP are separated from their residential neighbors,
the resulting districts should have \emph{lower} within-district JPR variance
than the baseline (districts contain either mostly residential tracts or mostly
employment-center tracts, not a mix of both).
This prediction will be validated in Phase~2 along with the full 50-state
application.

\subsection{Cross-Algorithm Comparison}

Table~\ref{tab:algorithm-comparison} places M.1 in context with other weight
modes evaluated on NC-14.

\begin{table}[ht]
\centering
\caption{NC-14 proportionality gap across algorithm configurations.
  Economic character produces the largest improvement among single-seed
  standard-bisect Layer~2 modifications.}
\label{tab:algorithm-comparison}
\begin{tabular}{llr}
  \toprule
  Weight mode & Config & Gap (pp) \\
  \midrule
  Geographic (baseline) & \texttt{--weights-override geographic} & $-$20.7 \\
  County (T.3) & \texttt{--weights-override county} & varies \\
  Economic character (M.1/M.9) & \texttt{--weights-override economic-character} & $-$6.4 \\
  \bottomrule
\end{tabular}
\end{table}

The economic character result of $-$6.4\,pp is the closest to proportionality
of any Layer~2 modification on NC-14 using standard-bisect structure.
Note that the AR algorithm (T.4, ApportionRegions) and convergence search
(U.1) can produce further improvements; M.1 is a Layer~2 modification and
composes with those Layer~1 and Layer~3 choices.
